\chapter{Programming models: OpenACC and CUDA Fortran}\label{ch:gpu}

A GPU programming model answers three questions:
\begin{itemize}
  \item which code runs on the device;
  \item where the data lives, and when it moves;
  \item how the iterations of a loop are mapped onto threads.
\end{itemize}
The port writes every kernel in OpenACC \citep{openacc33}, a set of directives in the existing Fortran, and uses
CUDA Fortran \citep{cudafortran} only for kernels that profiling shows to be the bottleneck. Both are compiled by
NVIDIA's \code{nvfortran}. This chapter explains both, with the constructs the port uses, and the reasons for the
choice.

\section{OpenACC}

\subsection{Execution model}

OpenACC describes parallelism at three levels:
\begin{itemize}
  \item \emph{gangs}, which run independently of each other;
  \item \emph{workers} within a gang;
  \item \emph{vector lanes} within a worker, which execute in lockstep.
\end{itemize}
On an NVIDIA GPU, \code{nvfortran} maps a gang to a thread block and the vector lanes to the threads of the block,
128 by default. Workers, if used, become the warps of a block. The port uses only gangs and vector lanes: every
kernel is a \code{parallel loop gang vector}, whose collapsed iteration space is split into blocks of
\code{vector_length} threads.

\subsection{The kernel directive}

Every kernel of the port has the same form (\code{port/agent/CODING_STANDARD.md}, section~4):
\begin{lstlisting}
!$acc parallel loop gang vector collapse(N) if(gpu_on(R_X)) default(none) &
!$acc& present(<arrays>) firstprivate(<scalars read>) private(<scalars written>)
\end{lstlisting}
Each clause has a purpose:
\begin{description}
  \item[\code{collapse(N)}] Flattens $N$ perfectly nested loops into one iteration space: \code{collapse(3)} over
    $(j,k,i)$ for pointwise kernels, \code{collapse(2)} over $(j,i)$ for column kernels.
  \item[\code{if(gpu_on(R_X))}] When the condition is false, the region runs on the host over host data. This is
    OpenACC's \emph{host fallback}, and it is what lets one routine be switched between device and host at run time
    (\cref{ch:port}).
  \item[\code{default(none)}] Every variable used in the region must be listed in a clause. The compiler reports what
    is missing at compile time, so a forgotten \code{private}, which would be a data race, cannot pass unnoticed.
  \item[\code{present(...)}] The arrays must already be on the device; if one is not, the run stops. A data clause
    such as \code{copyin} would instead copy the array at every launch: correct, but hundreds of times slower.
  \item[\code{firstprivate}, \code{private}] Scalars read in the loop get their host value in each thread; scalars
    written in the loop get a private copy per thread.
\end{description}

Inner loops that must run sequentially inside a thread, such as the $k$ loop of a column or a loop over species, are
marked \texttt{!\$acc loop seq}. They then run in the source order and direction. (A \code{DO WHILE} loop, or a loop with
an \code{EXIT}, takes no directive at all: OpenACC forbids one there, and it runs sequentially anyway.)

\Cref{lst:calc-alt-acc} shows a complete pointwise kernel: \code{calc_alt}, which forms the full inverse density
$\alpha = \alpha' + \bar\alpha$ (\cref{ch:eq}).

\begin{figure}[ht]
\begin{lstlisting}
      itf=MIN(ite,ide-1)
      jtf=MIN(jte,jde-1)
      ktf=MIN(kte,kde-1)
! K-PREP-7: full inverse density (Template A)
!$acc parallel loop gang vector collapse(3) if(gpu_on(R_CALC_ALT)) default(none) &
!$acc& present(alt, al, alb) firstprivate(its, itf, jts, jtf, kts, ktf)
      DO j=jts,jtf
      DO k=kts,ktf
      DO i=its,itf
        alt(i,k,j) = al(i,k,j)+alb(i,k,j)
      ENDDO
      ENDDO
      ENDDO
\end{lstlisting}
\caption{\code{calc_alt} as an OpenACC kernel (\code{dyn_em/module_big_step_utilities_em.F}; the island around it is
omitted, \cref{ch:port}). The loops are the CPU loops, unchanged.}
\label{lst:calc-alt-acc}
\end{figure}

\subsection{Data}

OpenACC assumes that host and device have separate memories, and that each array may have a copy on both. The port
uses these constructs:
\begin{center}
\begin{tabular}{ll}
\toprule
What & OpenACC \\
\midrule
allocate on the device and copy, once & \texttt{!\$acc enter data copyin(a)} or \code{create(a)} \\
release & \texttt{!\$acc exit data delete(a)} \\
copy host to device & \texttt{!\$acc update device(a)} \\
copy device to host & \texttt{!\$acc update self(a)} \\
module variable on the device & \texttt{!\$acc declare create(t)} in the module \\
is it there? & \code{acc_is_present(a)} \\
\bottomrule
\end{tabular}
\end{center}
The model's
state is copied to the device once, when each domain is allocated, and stays there (\cref{ch:port}). All other
movement is an explicit \code{update} at a place the port can name.

\subsection{Procedures}

A procedure called inside a kernel must be compiled for the device too. \texttt{!\$acc routine seq} after its
\code{SUBROUTINE} or \code{FUNCTION} statement makes the compiler build a device version that runs sequentially in
the calling thread, beside the normal host version. The port uses it for every column-physics scheme, for
\code{fire_ros}, and for the reproducible math functions.

\subsection{Asynchrony}

By default, a kernel launch waits until the kernel has finished. With \code{async(q)} the host continues at once and
the kernel joins queue $q$; kernels in one queue run in order, kernels in different queues may overlap, and
\texttt{!\$acc wait(q)} waits for a queue. On NVIDIA GPUs a queue is a CUDA stream. The port leaves asynchrony to the
performance stage (\cref{ch:fusion}), because it changes only the schedule, never the results.

\section{CUDA Fortran}

CUDA Fortran is NVIDIA's extension of Fortran with the concepts of CUDA \citep{cudafortran}:
\begin{itemize}
  \item \code{attributes(global)} marks a kernel, \code{attributes(device)} a routine called from kernels;
  \item arrays with the \code{device} attribute live in device memory;
  \item a kernel is launched with an explicit grid and block shape,
    \texttt{call k<{}<{}<grid, block, shmem, stream>{}>{}>(args)};
  \item inside a kernel, \code{blockIdx}, \code{blockDim} and \code{threadIdx} give the thread's position;
  \item \code{attributes(shared)} declares shared memory, and \code{call syncthreads()} synchronizes a block;
  \item warp-level functions exchange values between the threads of a warp (\code{__shfl_down}) or vote on a
    condition (\code{ballot}, \code{allthreads}, \code{anythread}).
\end{itemize}
The programmer chooses everything that OpenACC leaves to the compiler: the shape of the blocks, what each thread
does, what is staged in shared memory. \Cref{lst:calc-alt-cuf} shows \code{calc_alt} in this form. It is an
illustration; the port writes CUDA Fortran only where a measurement justifies it.

\begin{figure}[ht]
\begin{lstlisting}
attributes(global) subroutine calc_alt_cuf(alt, al, alb, ims, ime, kms, kme, jms, jme, &
                                           its, itf, kts, ktf, jts, jtf)
   integer, value :: ims, ime, kms, kme, jms, jme, its, itf, kts, ktf, jts, jtf
   real, device :: alt(ims:ime,kms:kme,jms:jme), al(ims:ime,kms:kme,jms:jme), &
                   alb(ims:ime,kms:kme,jms:jme)
   integer :: i, k, j
   i = its + (blockIdx%x-1)*blockDim%x + threadIdx%x - 1
   k = kts + blockIdx%y - 1
   j = jts + blockIdx%z - 1
   if (i <= itf) alt(i,k,j) = al(i,k,j)+alb(i,k,j)
end subroutine calc_alt_cuf

! launch: one block of 128 threads per 128 values of i, one block row per (k,j)
!$acc host_data use_device(alt, al, alb)
call calc_alt_cuf<<<dim3((itf-its+128)/128, ktf-kts+1, jtf-jts+1), dim3(128,1,1)>>>( &
     alt, al, alb, ims, ime, kms, kme, jms, jme, its, itf, kts, ktf, jts, jtf)
!$acc end host_data
\end{lstlisting}
\caption{\code{calc_alt} in CUDA Fortran (illustration). The statement is the same; the mapping of iterations to
threads is written out by hand. \code{host_data use_device} hands the device addresses of OpenACC's copies to the
CUDA Fortran kernel.}
\label{lst:calc-alt-cuf}
\end{figure}

\paragraph{Working together.} The two models share one compiler and one device memory:
\begin{itemize}
  \item \texttt{!\$acc host\_data use\_device(a)} passes the device address of an OpenACC array to a CUDA Fortran
    kernel, so the state never needs a second copy;
  \item \code{acc_get_cuda_stream(q)} gives the CUDA stream of OpenACC queue $q$, so a CUDA Fortran kernel can run in
    order with the OpenACC kernels around it;
  \item the build flag \code{-cuda} enables CUDA Fortran beside \code{-acc=gpu}.
\end{itemize}

\begin{portnote}
A CUDA Fortran kernel lives next to its OpenACC version, under \texttt{\#ifdef WRF\_CUF\_<KERNEL>}, and must give the same
bits: the same statements in the same order (\code{CODING_STANDARD.md} section~11). The OpenACC version stays the
reference. The test runs both on the same window and compares the traces.
\end{portnote}

\section{What the compiler tells}

\code{nvfortran -Minfo=accel} reports, for every compute region:
\begin{itemize}
  \item whether a device kernel was generated;
  \item how each loop was scheduled (gang, vector, sequential);
  \item every data movement the compiler added implicitly.
\end{itemize}
The port requires that every kernel is offloaded, that its inner loops are sequential where the source has a
recurrence, and that no implicit data movement appears (rung L4 of the routine ladder, \cref{ch:port}). With
\code{-gpu=ptxinfo} the compiler also reports the registers, spills and stack frame of each kernel: the first layer of
the profiler (\cref{ch:prof}).

The flags that keep the arithmetic reproducible belong to the same command line: \code{-gpu=nofma,noflushz} for the
device and \code{-Mnofma -Mnoflushz -Kieee} for the host (\cref{ch:repro,ch:iobuild}).

\section{Why this pair}

The choice was made twice (\code{roadmap/decisions/ADR-001-gpu-programming-model.md}).

\paragraph{Revision 1: OpenMP offload.} The first decision chose OpenMP~5 \code{target} directives
\citep{openmp52}, because OpenMP offload is the one standard that every GPU vendor supports in its own Fortran
compiler. The port's first kernels and its compiler probes were written in OpenMP.

\paragraph{Revision 2: OpenACC with CUDA Fortran.} The project's owner then chose performance over portability: the
point of the port is speed on the GPUs the project has, NVIDIA A100 and H100. On NVIDIA hardware OpenACC is the more
mature model in \code{nvfortran}, and CUDA Fortran gives full control where it is needed, without leaving Fortran.
The kernels written before the switch were converted mechanically (\code{port/tools/omp2acc.py}); the kernel form,
the tests and the bitwise method did not change.

What the pair keeps:
\begin{itemize}
  \item \textbf{One Fortran source.} The kernels are the CPU loops with directives in front. The line-by-line
    correspondence with NCAR's source, on which both the bitwise verification and the forward port to newer WRF
    versions rely (\cref{ch:versions}), is preserved.
  \item \textbf{Checks without a GPU.} gfortran with \code{-fopenacc} compiles the GPU view of the source and runs its
    OpenACC regions on the host. The arithmetic of every restructured kernel can therefore be checked on any CPU, and
    \code{default(none)} omissions are reported by two compilers.
  \item \textbf{Control of the arithmetic.} No FMA contraction, no flush to zero, and the portable elementary
    functions, on both host and device (\cref{ch:repro}).
\end{itemize}
What it gives up: AMD and Intel GPUs. OpenACC reaches AMD GPUs only through the HPE Cray and GCC compilers, and CUDA
Fortran not at all. A later port to AMD would rewrite the directives as OpenMP offload (the reverse of
\code{omp2acc.py}) and translate the few CUDA Fortran kernels to HIP. The bitwise method would carry over unchanged:
contraction off, the portable math library, and the same probes and tests.

\begin{portnote}
Every construct the port uses is first tried in a small probe program on the pinned compiler version: \code{if}
host fallback, presence of arrays mapped through pointers, \code{declare create} on module allocatables, calls of
\code{routine seq} procedures, private arrays, reductions, the device stack size. The probes of revision 1 are in
\code{port/tests/omp_features}; their OpenACC versions are phase S0-05 of the plan, and must pass before any kernel
is believed.
\end{portnote}
